\section{Background and Motivation}
\subsection{Inefficiencies in remote meetings: Challenges in convergent decision-making processes}
Convergent decision-making in remote meetings faces unique obstacles:
\begin{itemize}
  \item Off-topic discussions: Studies show that 30–50\% of meeting time is spent on irrelevant tangents,
    derailing progress [Rogelberg et al., 2006]. \cite{chenAreWeTrack2025}
  \item Unequal participation: Dominant voices can overshadow quieter team members, leading to groupthink and
    suboptimal outcomes [Nunamaker et al., 1991].
  \item Lack of structure: Without clear facilitation, meetings often lack timeboxing, role assignment, or
    explicit decision criteria, reducing efficiency [Schwartz, 2016].
\end{itemize}
Traditional solutions—such as hiring professional facilitators or training team members—are cost-prohibitive
or scalability-limited. For example, a certified facilitator may charge €100–300/hour [International
Association of Facilitators, 2024], making continuous support impractical for most organizations.
Furthermore, not every team has permanent access to trained personnel, and ad-hoc facilitation often fails to
address systemic issues like dysfunctional team dynamics or blind spots in group perception [Edmondson, 1999].

\subsection{Potential of AI-driven facilitation to enhance goal orientation }

AI can bridge this gap by:
\begin{itemize}
  \item Automating real-time analysis: Natural Language Processing (NLP) can transcribe and semantically
    analyze discussions to identify deviations from the meeting’s goal [Bommasani et al., 2021].
  \item Providing adaptive interventions: Context-sensitive prompts (e.g., "Let’s return to the agenda:
    prioritizing the top 3 solutions") can refocus conversations without human intervention.
  \item Democratizing participation: AI can nudge underrepresented voices (e.g., "Alex, you haven’t shared
    your perspective yet—would you like to add anything?") or summarize key points to ensure alignment.
\end{itemize}
Pilot studies on AI moderation tools (e.g., Humley’s AI meeting assistant) report 20–30\% reductions in
meeting duration and improved participant satisfaction [Humley, 2023]. However, these tools often lack
proactivity—reacting to issues rather than preventing them—and customization for specific use cases like
convergent decision-making.
